% answers_short.tex —— 2004 年·普通高等学校招生考试·湖南卷·理科 · 简短答案速查
% 由 tools/make_short_answers.py 自 answers.tex 的 \gksolhead/\gkanswer 提取，
% 并与 2004/2004_数学真题_湖南理科_答案要点.md 逐题比对后生成，勿手改（改源 MD 后重跑）。

\begin{gknotice}
\textbf{使用说明}：本节只给最终答案便于快速核对，完整推导见解析版。答案为编写者独立推导，\textbf{非官方标准答案}。
\end{gknotice}

\gkpart{选择题}
\begin{center}\begin{tabular}{ccccccccccccc}
\toprule
题号 & 1 & 2 & 3 & 4 & 5 & 6 & 7 & 8 & 9 & 10 & 11 & 12 \\
\midrule
答案 & D & A & B & C & B & C & B & C & A & C & B & D \\
\bottomrule
\end{tabular}\end{center}

\gkpart{填空题}
\begin{itemize}[leftmargin=2.6em]
  \item[\textbf{13.}] $4$
  \item[\textbf{14.}] $\dfrac{3}{4}$
  \item[\textbf{15.}] $9$
  \item[\textbf{16.}] $\left(0, \dfrac{2}{11}\right]$
\end{itemize}

\gkpart{解答题}
\begin{itemize}[leftmargin=2.6em]
  \item[\textbf{17.}] $\dfrac{5\sqrt{3}}{2}$
  \item[\textbf{18.}] (1) 见解析版；(2) $30^\circ$（或 $\dfrac{\pi}{6}$）；(3) 在棱 $PC$ 上存在点 $F$，当点 $F$ 为棱 $PC$ 的中点时，$BF \parallel$ 平面 $AEC$，证明见解析
  \item[\textbf{19.}] (1) 甲、乙、丙加工零件是一等品的概率分别为 $\dfrac{1}{3}$，$\dfrac{1}{4}$，$\dfrac{2}{3}$；(2) $\dfrac{5}{6}$
  \item[\textbf{20.}] (1) 当 $a = 0$ 时，单调递减区间为 $(-\infty, 0]$，单调递增区间为 $[0, +\infty)$；当 $a < 0$ 时，单调递减区间为 $(-\infty, 0]$ 和 $\left[-\dfrac{2}{a}, +\infty\right)$，单调递增区间为 $\left[0, -\dfrac{2}{a}\right]$；(2) 当 $-2 \leqslant a \leqslant 0$ 时，最大值为 $\mathrm{e}^a$；当 $a < -2$ 时，最大值为 $\dfrac{4}{a^2 \mathrm{e}^2}$
  \item[\textbf{21.}] (1) 见解析版；(2) $\left(x + \dfrac{3}{2}\right)^2 + \left(y - \dfrac{23}{2}\right)^2 = \dfrac{125}{2}$（或 $x^2 + y^2 + 3x - 23y + 72 = 0$）
  \item[\textbf{22.}] (1) 见解析版；(2) $x_n = 1 - 2\left(\dfrac{1}{2k}\right)^n$；(3) 当 $0 < |k| < \dfrac{1}{2}$ 时，对一切 $n \in \mathbf{N}^*$ 恒有 $2|PP_n|^2 > 4k^2|PP_1|^2 + 5$；当 $|k| > \dfrac{1}{2}$ 时，对一切 $n \in \mathbf{N}^*$ 恒有 $2|PP_n|^2 < 4k^2|PP_1|^2 + 5$
\end{itemize}
